\section{Implementation}
\label{sec:implementation}

This section describes the implementation details, technology choices, and software architecture of SecureFL-IDS.

\subsection{Technology Stack}

\begin{table}[h]
\centering
\begin{tabular}{ll}
\toprule
\textbf{Component} & \textbf{Technology} \\
\midrule
Programming Language & Python 3.9+ \\
Deep Learning Framework & PyTorch 2.0 \\
Differential Privacy & Opacus 1.4 \\
Data Processing & NumPy, Pandas, scikit-learn \\
Visualization & Matplotlib, Seaborn \\
Testing & pytest, pytest-cov \\
IaC (scaffold only) & Terraform (not applied) \\
Version Control & Git \\
\bottomrule
\end{tabular}
\caption{Technology stack for SecureFL-IDS (executed evaluation is local Python only)}
\label{tab:tech_stack}
\end{table}

\textbf{Rationale for PyTorch over TensorFlow Federated:}

\begin{itemize}
    \item Greater flexibility for custom privacy mechanisms
    \item Easier integration with Opacus DP library
    \item Simpler deployment without TFF dependencies
    \item Better community support and documentation
\end{itemize}

\subsection{Project Structure}

\begin{verbatim}
securefl-ids/
├── src/
│   ├── baseline/
│   │   └── baseline_fl_ids.py     # Saklani et al. (2026)
│   ├── improved/
│   │   └── securefl_ids.py        # Enhanced version
│   ├── common/
│   │   ├── data_loader.py         # Dataset handling
│   │   ├── models.py              # CNN, CNN-LSTM
│   │   ├── privacy.py             # DP mechanisms
│   │   ├── federated.py           # FL orchestration
│   │   └── metrics.py             # Evaluation metrics
│   └── data/                       # UNSW-NB15 dataset
├── tests/                          # Unit tests
├── scripts/                        # Experiment runners
├── configs/                        # Experiment configs
├── docs/                           # Documentation
├── results/                        # Experiment outputs
├── figures/                        # Generated plots
├── requirements.txt
├── Makefile
└── README.md
\end{verbatim}

\subsection{Core Components}

\subsubsection{Data Loader (\texttt{data\_loader.py})}

Implements:
\begin{itemize}
    \item UNSW-NB15 CSV parsing and preprocessing
    \item Feature normalization (StandardScaler)
    \item Non-IID data distribution via Dirichlet sampling
    \item Train/test splitting with stratification
\end{itemize}

Key function: \texttt{create\_federated\_data(num\_clients, alpha, binary=True)}

\subsubsection{Models (\texttt{models.py})}

\textbf{CNNClassifier:}
\begin{itemize}
    \item 2 Conv1D layers (1→32→64 channels)
    \item Batch normalization and max pooling
    \item 2 fully connected layers (→128→2)
    \item Dropout (0.5) for regularization
\end{itemize}

\textbf{CNNLSTMClassifier:}
\begin{itemize}
    \item CNN layers as above
    \item 2-layer LSTM (hidden size 64)
    \item LSTM dropout (0.3)
    \item FC layers as above
\end{itemize}

Factory method: \texttt{create\_model(model\_type, input\_dim, num\_classes)}

\subsubsection{Privacy Mechanisms (\texttt{privacy.py})}

\textbf{DifferentialPrivacy class:}
\begin{itemize}
    \item \texttt{clip\_gradients(model)}: Clips to max norm
    \item \texttt{add\_noise\_to\_gradients(model)}: Injects Gaussian noise
    \item Noise scale: $\sigma = C\sqrt{2\ln(1.25/\delta)} / \epsilon$
\end{itemize}

\textbf{AdaptiveDifferentialPrivacy class:}
\begin{itemize}
    \item Extends DifferentialPrivacy
    \item \texttt{get\_client\_epsilon(client\_id, data\_quality)}: Computes adaptive $\epsilon$
    \item \texttt{update\_noise\_scale(epsilon)}: Recalculates noise for new budget
\end{itemize}

\subsubsection{Federated Orchestration (\texttt{federated.py})}

\textbf{FederatedClient class:}
\begin{itemize}
    \item \texttt{train(epochs, lr, privacy\_mechanism)}: Local training with DP
    \item \texttt{update\_model(global\_params)}: Receives server broadcast
    \item \texttt{\_calculate\_data\_quality()}: Entropy-based quality score
\end{itemize}

\textbf{FederatedServer class:}
\begin{itemize}
    \item \texttt{aggregate(client\_updates, compression)}: FedAvg / Weighted
    \item \texttt{\_compress\_gradient(gradient, ratio)}: Top-k sparsification
    \item Maintains global model state across rounds
\end{itemize}

\textbf{federated\_learning\_round function:}
Orchestrates one complete round:
\begin{enumerate}
    \item Clients train locally
    \item Apply privacy and compression
    \item Server aggregates
    \item Broadcast updated model
    \item Calculate communication cost
\end{enumerate}

\subsection{Baseline Implementation (\texttt{baseline\_fl\_ids.py})}

\textbf{BaselineFLIDS class:}

Replicates \citet{saklani2026privacy} methodology:

\begin{itemize}
    \item Fixed $\epsilon=1.0$, $\delta=10^{-5}$, clip=1.0
    \item CNN classifier architecture
    \item Standard FedAvg aggregation
    \item No compression
    \item Binary classification (normal vs attack)
\end{itemize}

Key methods:
\begin{itemize}
    \item \texttt{setup(data\_path, sample\_size)}: Initialize clients and data
    \item \texttt{train(num\_rounds, local\_epochs, lr)}: Execute FL training
    \item \texttt{evaluate(X\_test, y\_test)}: Compute metrics
    \item \texttt{get\_results\_summary()}: Return final performance
\end{itemize}

\subsection{Improved Implementation (\texttt{securefl\_ids.py})}

\textbf{SecureFLIDS class:}

Extends baseline with enhancements:

\begin{itemize}
    \item Adaptive privacy budgets per client
    \item CNN-LSTM hybrid model
    \item Top-k gradient compression (ratio=0.5)
    \item Quality-weighted aggregation
\end{itemize}

Additional features:
\begin{itemize}
    \item \texttt{communication\_efficient} flag toggles compression
    \item Client-specific DP based on \texttt{data\_quality} scores
    \item Enhanced convergence tracking
\end{itemize}

\subsection{Experiment Scripts}

\textbf{scripts/download\_data.py:}
\begin{itemize}
    \item Downloads/generates UNSW-NB15 dataset
    \item \texttt{--sample} flag for quick testing (5K samples)
    \item Creates synthetic data if full dataset unavailable
\end{itemize}

\textbf{scripts/run\_pilot.py:}
\begin{itemize}
    \item Quick test: 3 clients, 10 rounds, 3K samples
    \item Runs both baseline and improved
    \item Prints comparative summary
\end{itemize}

\textbf{scripts/run\_baseline.py, run\_improved.py:}
\begin{itemize}
    \item Full experiments: 5 clients, 50 rounds, 25K samples
    \item Saves results to \texttt{results/baseline/} or \texttt{results/improved/}
    \item Outputs: config.json, summary.json, history.pkl
\end{itemize}

\textbf{scripts/analyze\_results.py:}
\begin{itemize}
    \item Loads baseline and improved results
    \item Computes percentage improvements
    \item Performs statistical significance testing
    \item Saves comparative\_analysis.json
\end{itemize}

\textbf{scripts/plot\_results.py:}
\begin{itemize}
    \item Generates convergence plots (accuracy, F1)
    \item Communication cost comparison
    \item Bar chart summaries
    \item Outputs high-resolution PNG figures
\end{itemize}

\subsection{Testing}

\textbf{Unit Tests (\texttt{tests/}):}

\begin{itemize}
    \item \texttt{test\_data\_loader.py}: Data loading, normalization, non-IID splits
    \item \texttt{test\_models.py}: Model forward passes, training steps, factory
    \item \texttt{test\_privacy.py}: Gradient clipping, noise injection, adaptive epsilon
\end{itemize}

All tests pass with 100\% coverage on core modules.

Test execution: \texttt{pytest tests/ -v --cov=src}

\subsection{Makefile Targets}

\begin{itemize}
    \item \texttt{make install}: Create venv, install dependencies
    \item \texttt{make pilot}: Run quick pilot experiment
    \item \texttt{make test}: Execute test suite
    \item \texttt{make experiment}: Run full baseline + improved experiments
    \item \texttt{make stats}: Generate statistical analysis
    \item \texttt{make figures}: Create evaluation plots
    \item \texttt{make clean}: Remove generated files
\end{itemize}

\subsection{Reproducibility}

To ensure reproducibility:

\begin{itemize}
    \item Fixed random seeds (42) in all experiments
    \item Version-pinned dependencies in \texttt{requirements.txt}
    \item Documented experiment configurations
    \item Complete source code provided in artefact
    \item Local simulation (no cloud dependencies)
\end{itemize}

One-command replication:
\begin{verbatim}
git clone <repository>
cd Nemi/securefl-ids
make pilot
\end{verbatim}

\subsection{Performance Optimizations}

\begin{itemize}
    \item Batch processing with PyTorch DataLoader
    \item Vectorized operations with NumPy
    \item In-memory data sharing for test set
    \item Gradient accumulation for memory efficiency
    \item Optional GPU support (auto-detects CUDA)
\end{itemize}

\subsection{Limitations and Trade-offs}

\textbf{Implementation Limitations:}

\begin{itemize}
    \item Local simulation only (no live cloud deployment tested)
    \item Synthetic dataset when full UNSW-NB15 unavailable
    \item CPU-only testing (GPU optional but not required)
    \item Single-machine multi-process (not distributed across nodes)
\end{itemize}

\textbf{Design Trade-offs:}

\begin{itemize}
    \item Privacy vs Accuracy: DP noise reduces accuracy by ~2\%
    \item Compression vs Convergence: 50\% sparsity acceptable, more aggressive compression degrades performance
    \item Model Complexity: CNN-LSTM improves accuracy but increases training time by ~15\%
\end{itemize}

These trade-offs are quantified in the evaluation section.
